\section{The unresolved step and the branch removed here}
The universal assertion remains that every prescribed hard prime supplies an
original positive exterior or middle state. The complete atlases from the
preceding work decide a fixed input; their finite index bounds do not prove
occupancy. We do not advance to a universally successful algorithm under an
assumed occupancy statement.

The previous exterior cutoff has an asymptotically sharp integer family. Here
we determine its exact marked cubic geometry, prove that its diagonal cofactor
locus already has an original middle return, and classify every hard-prime
exterior state whose two shape defects have absolute value at most seven.
Every such state is on that same singular family, and has two original middle
orientations at its \emph{unchanged} distinguished denominator. At radius eight
there is at most one nonsingular shape; its grade congruence also gives an explicit
middle return, independently of solving its associated quartic.
This is an existence-preserving branch reduction, not a proof of the remaining
unbounded positive coefficient.

The original Type-I/II parametrizations are antecedents \cite[\S2]{ET}.
The integer family itself occurs in the preceding exterior proof
\cite[\S5]{Exterior}. No historical-priority determination is made for the
present deductions. All algebra and the finite modular classification used here
are proved below; no theorem about integral points on elliptic curves is imported.

\section{Original source and the marked cubic}
Let $p\equiv1\pmod8$ be prime. An original exterior state has
\[
 p/4<a<p/2,\quad R=4a-p,\quad u\mid a^2,\quad R\mid4u+1.
\]
Retain $a=\prod_l l^{e_l}$, $\beta_l=v_l(u)-e_l$, and the ordered normalization
\[
 g=\gcd(a,u),\qquad
 (h,r,s)=(g^2/u,u/g,a/g),\qquad \kappa=(pr+s)/R.
\]
Then $a=hrs$, $u=hr^2$, $\gcd(r,s)=1$, and $r<\kappa$. The ordered return is
\[
 (x,y,z)=(a,hs\kappa,phr\kappa),\qquad
 \frac4p=\frac1x+\frac1y+\frac1z.                         \tag{1}
\]
Its ordered inverse is $u=a^2/(Ry-pa)$. Prime valuations prove the integrality
of the normalization. Multiplication proves (1), without identifying its
ordered denominators with another channel's order.

Keep the original complement and cofactor
\[
 v=a^2/u=hs^2,\qquad D=(4u+1)/R=(r+\kappa)/s.
\]
They satisfy
\[
 RD=4u+1,\quad pD+1=4hr\kappa,\quad
 (p+R)^2+4v=4vRD,\quad R,D\equiv3\pmod4.                 \tag{2}
\]
The case $h=1$ would make $-1$ a square modulo $R\equiv3\pmod4$. Thus $h>1$.
Also $\gcd(h,R)=\gcd(h,D)=1$, and consequently $v\ne R,D$.
No assertion $\gcd(v,D)=1$ is made.

Define the \emph{marked shape} and its coefficient
\[
 \alpha=v-R,\quad \beta=v-D,\quad B=\alpha\beta-1,
 \qquad X=v,\quad Y=2a.                                  \tag{3}
\]
(The symbol $\beta$ in (3) is not an original exponent $\beta_l$.)
The pairs $(\alpha,\beta)$ and $(R,D)$ are ordered data: interchanging them
leaves the unmarked cubic equation unchanged but changes the reconstructed
distinguished residual and, in general, the reconstructed prime.
\begin{theorem}[Exact ordered shape curve and inverse]\label{thm:curve}
Every original state maps to an integral point
\[
 \boxed{Y^2=X^3-(\alpha+\beta)X^2+BX.}                    \tag{4}
\]
The complete inverse is
\[
 R=X-\alpha,\quad D=X-\beta,\quad a=Y/2,\quad
 p=2Y-R,\quad u=(RD-1)/4.                                \tag{5}
\]
Its domain consists exactly of points of (4) with $X>0$, positive even $Y$,
positive $R,D\equiv3\pmod4$, $Y>R$, and the prime condition on the reconstructed
$p\equiv1\pmod8$. The original normalization then reconstructs every marking.
Moreover
\[
 \boxed{h\mid B,\qquad p\equiv\alpha\pmod h,\qquad
 p\beta\equiv1\pmod h.}                                 \tag{6}
\]
The Weierstrass discriminant of (4) is
\[
 \boxed{16B^2((\alpha-\beta)^2+4).}                       \tag{7}
\]
Thus its only singular integer shapes are $(\alpha,\beta)=(1,1)$ and $(-1,-1)$.
\end{theorem}
\begin{proof}
From $a^2=uv$ and $4u=RD-1$,
$Y^2=v(RD-1)=v((v-\alpha)(v-\beta)-1)$, proving (4).
Conversely (4) gives $a^2=X(RD-1)/4=Xu$, so $u$ is a positive divisor of $a^2$.
The residue requirements make $u$ integral. Equation (5) gives $R=4a-p>0$ and
$2a<p$ precisely from $Y>R$. The original gate follows from $4u+1=RD$.
The prime condition and $a<p$ imply the original unit condition.
All formulas are literal inverses.

Reducing $RD-1=4u$ modulo $h$ gives $\alpha\beta-1\equiv0$; reducing
$p=4hrs-R$ and $pD+1=4hr\kappa$ gives the other two identities in (6).
For the cubic $X(X^2-(\alpha+\beta)X+B)$, the polynomial discriminant is
$B^2((\alpha+\beta)^2-4B)$. The Weierstrass discriminant is sixteen times
the polynomial discriminant, giving (7).
Its second factor is $(\alpha-\beta)^2+4>0$. Thus singularity means
$\alpha\beta=1$, whose integer possibilities are the two stated shapes.
\end{proof}

For a fixed original $p$ and shape, the point must also retain the line
\[
 \boxed{2Y=p+X-\alpha.}                                  \tag{8}
\]
Eliminating $Y$ yields the explicit cubic
\[
 4X^3-[4(\alpha+\beta)+1]X^2+
 [4B-2(p-\alpha)]X-(p-\alpha)^2=0.                        \tag{9}
\]
There are at most three original states at a fixed $(p,\alpha,\beta)$.
This is a fibre bound, not an existence assertion. The integer coefficient
pushforward that forgets $X,Y$ sums those actual source coordinates; its kernel
is generated by their same-shape differences. Retaining $X,Y$ gives the
bijection of Theorem~\ref{thm:curve}. The point $(0,0)$ of a nonsingular curve is
outside the original positive source.

\section{Original character constraints}
Here and below, a \emph{hard prime} means
\[
 p\bmod840\in\{1,121,169,289,361,529\}.
\]
\begin{lemma}\label{lem:char}
For every original square divisor, $(u/R)=(u/p)$, where the first symbol is
Jacobi. In particular $(h/p)=-1$. If $p$ is hard, every prime factor among
$2,3,5,7$ has character $+1$, and $h$ contains a prime factor greater than seven.
Neither $\alpha$ nor $\beta$ can be a positive integer square.
\end{lemma}
\begin{proof}
For an odd prime $l\mid a$, reciprocity and $R\equiv-p\pmod l$ give
$(l/R)=(-1/l)(-p/l)=(p/l)=(l/p)$. For 2 use $R\equiv-p\pmod8$ when $2\mid a$.
Multiply over the original exponents. The exterior gate has $u\equiv-1/4\pmod R$
and $R\equiv3\pmod4$, so $(u/p)=-1$; $u=hr^2$ gives the assertion for $h$.
For the six hard classes, $p\equiv1\pmod{24}$, $p\equiv1,4\pmod5$, and
$p\equiv1,2,4\pmod7$, proving the four residue assertions.
These are the ordinary reciprocity and supplementary laws \cite[\S6.4]{Hatcher}.

If $\alpha=c^2>0$, then every odd $l\mid h$ sees $p\equiv c^2\pmod l$ by (6).
The value is a nonzero square because $p\nmid h$. Thus $(l/p)=1$ for every such
$l$, while $(2/p)=1$, contradicting $(h/p)=-1$. If $\beta=c^2>0$, use
$p\beta\equiv1\pmod h$ instead. The shapes $\alpha=0$ or $\beta=0$ were already
excluded by $v\ne R,D$.
\end{proof}

\section{Every diagonal exterior source has an actual middle return}
\begin{theorem}[Cofactor collision crosses the channels]\label{thm:collision}
Suppose an original exterior state at $p\equiv1\pmod8$ has $R=D$.
Set
\[
 z=(R-1)/2=\delta t^2,
\]
where $\delta$ is the positive squarefree part (odd exponents retained).
Then $\delta\equiv3\pmod4$, $t$ is odd,
\[
 \delta\mid a,\quad \delta\mid p+1,\quad 3\le\delta<R.
\]
With $c=(p+1)/\delta$ and $j=(\delta+1)/4$, the explicit original middle state is
\[
 \boxed{(h_M,r_M,s_M,\lambda_M)=(j,1,c,1),\quad
 a_M=jc,\ R_M=c+1,\ u_M=j,\ Q_M=\delta.}                 \tag{10}
\]
It returns $(a_M,pa_M,pj)$, including the original complementary middle
orientation. Its small cofactor satisfies $\min(R_M,Q_M)<R$.
There is also the original endpoint exterior state
\[
 R_0=\delta,\quad a_0=u_0=(p+\delta)/4,\quad
 (h_0,r_0,s_0,\kappa_0)=(a_0,1,1,c),                    \tag{11}
\]
with $a_0<a$.
\end{theorem}
\begin{proof}
Since $R^2=4u+1$, $u$ is even, hence $a=hrs$ is even. Thus $R\equiv7\pmod8$
and $z\equiv3\pmod4$. The equality $u=z(z+1)$ shows that every prime factor of
$z$ divides $a$. Consequently its squarefree part $\delta$ divides $a$, and
$p=4a-R\equiv-1\pmod\delta$. Also $z=\delta t^2$ gives $t$ odd and
$\delta\equiv3\pmod4$.

Now $p=\delta c-1=4jc-c-1$. Thus (10) has $R_M=4a_M-p=c+1$,
$r_M+s_M=R_M\lambda_M$, and $u_M\mid a_M^2$.
The first-half upper bound is $c(\delta-1)>2$, true because $c\ge2$ and
$\delta\ge3$. All denominators are positive and multiplication verifies the
identity. Its second cofactor is exactly $\delta\le(R-1)/2$.
The endpoint in (11) follows by the same direct substitution; $\delta<R$
gives $a_0<a$. The endpoint formulas themselves are elementary antecedents;
the new inference here extracts their needed divisor from the \emph{actual}
diagonal source rather than assuming its existence.
\end{proof}

The map is not claimed injective after its original source is forgotten.
Its complete fibre above a marked (10) retains the positive odd integer $t$, with
\[
 R=2\delta t^2+1<p,\quad a=(p+R)/4,\quad
 u=\delta t^2(\delta t^2+1)\mid a^2.                      \tag{12}
\]
These explicit conditions and the squarefree marking are sufficient and necessary.
They reconstruct the source by gcd normalization. The finite bound on its fibre is
$t^2\le(p-3)/(2\delta)$. Source aggregation sums these integer fibres; it does
not replace them by a normalized count of residue classes.

The original residual itself need not decrease under the middle return.
At $p=2129$, the exterior source $(a,u,R,D)=(570,5700,151,151)$ has
$z=75=3\cdot5^2$. Its middle return has $Q_M=3$ but $R_M=711$.
The strict measure in the theorem is the smaller cofactor, not $R_M$ alone.
At the hard prime $p=48049$, $(a,u,R,D)=(12090,24180,311,311)$ gives
$\delta=155,t=1,j=39,c=310$ and the actual middle identity
\[
 \boxed{\frac4{48049}=\frac1{12090}+\frac1{580912410}
                         +\frac1{1873911}.}              \tag{13}
\]
The old E divisor $24180$ is not an M divisor at that shell; the return uses $39$.

\section{Complete hard-prime rigidity at shape radius seven}
\begin{theorem}[Small-shape rigidity]\label{thm:rigid}
For a hard prime and an original E state with
\[
 \max(|v-R|,|v-D|)\le7,
\]
there is an even integer $n\ge2$ such that
\[
 \boxed{h=4n^2-2,\quad r=n,\quad s=1,\quad
 R=D=4n^2-1,\quad\kappa=4n^2-n-1,}                       \tag{14}
\]
\[
 p=16n^3-4n^2-8n+1,\quad a=hn,\quad u=hn^2.
\]
At the \emph{same} original shell it forces the two middle states
\[
 \boxed{u_M=n,\quad u_M'=nh^2,\quad
 (h_M,r_M,s_M,\lambda_M)=(n,1,h,1)}                       \tag{15}
\]
and the complementary orientation has
\[
 (h_M',r_M',s_M',\lambda_M')=(n,h,1,1),\qquad
 Q_M'=4a-1.
\]
The first new ordered denominator triple is
$(nh,pnh,pn)$, and $Q_M=4n-1<R$.
\end{theorem}
\begin{proof}
There are 52 ordered coefficient profiles with $\alpha,\beta\in[-7,7]\setminus\{0\}$
and $\alpha\equiv\beta\pmod4$: the four residue groups have sizes $2,4,4,4$.
Their complete $B=\alpha\beta-1$ spectrum, written as value followed by
multiplicity, is
\[
\begin{array}{c|rrrrrrrrrrrrrrrrrr}
B&-37&-36&-17&-16&-13&-8&-5&-4&0&3&4&8&11&15&20&24&35&48\\
\#&2&4&2&4&4&4&2&4&2&2&4&2&4&2&4&2&2&2.
\end{array}
\]
If $B\ne0$, direct multiplication of these small integers gives the following
complete list of profiles for which $|B|$ has a prime factor greater than seven:
\[
\begin{array}{c|c|c}
 B& (\alpha,\beta)&h\text{ forced by }h\mid B,\ (h/p)=-1\\ \hline
 11&(2,6),(6,2),(-2,-6),(-6,-2)&11\\
 -13&(-6,2),(2,-6),(-2,6),(6,-2)&13\\
 -17&(-4,4),(4,-4)&17\\
 -37&(-6,6),(6,-6)&37.
\end{array}                                             \tag{16}
\]
Every other nonzero $B$ is supported on $2,3,5,7$, contrary to
Lemma~\ref{lem:char}. No unbounded factor search is implicit in (16).
Substituting $v=hs^2$ and $Y=2hrs$ into (4) gives
\[
 \boxed{4r^2=hs^4-(\alpha+\beta)s^2+B/h.}                \tag{17}
\]

For $h=11$, $\alpha\equiv2\pmod4$ would require $11s^2\equiv1\pmod4$ from
$R\equiv3$, impossible. For $h=17$, $\alpha\equiv0$ would require
$17s^2\equiv3\pmod4$, also impossible.

For $h=13$, the same residual condition forces $s$ odd. If $\alpha+\beta=4$,
(17) is $4r^2\equiv8\pmod{16}$, impossible. If $\alpha+\beta=-4$, it gives
$r$ even. Then $v\equiv5\pmod8$ and $\alpha\equiv2\pmod8$, so $R\equiv3$
and $p=4hrs-R\equiv5\pmod8$, again impossible.

For $h=37$, $s$ is odd and (17) makes $r$ odd. The condition $p\equiv1\pmod8$
selects $\alpha=-6,\beta=6$. Modulo three, (17) excludes $3\mid s$ and forces
$3\mid r$. But then $p=4hrs-37s^2-6\equiv2\pmod3$, contrary to the hard class.
This excludes every nonsingular profile.

When $B=0$, the positive-square lemma excludes $(1,1)$, leaving $(-1,-1)$.
Here $v=R-1=D-1$, and $RD-1=4hr^2$ becomes
\[
 s^2(v+2)=4r^2.
\]
Coprimality implies $s\mid2$. The choice $s=2$ makes $r$ odd and would give
$v+2=r^2$ with $v=4h$, impossible modulo four. Hence $s=1$, $h=v=4r^2-2$,
which proves (14). Reduction of $p$ modulo eight forces $r=n$ even.

For the same-shell return, $n+a=n(h+1)=nR$. Thus $u_M=n\mid a^2$ is an
original middle hit. Gcd normalization gives (15) and its complement. These
are two distinct oriented middle states, not one E state with a new label.
\end{proof}

The map from the singular E subfamily to the M orientation $r_M=1$ is bijective
onto the subfamily in (15): read $n=h_M$, then $h=4n^2-2$ and every original
coordinate in (14). The two M orientations retain one additional Boolean.
The corresponding endpoint E return has the particularly simple formulas
\[
 R_0=4n-1,\quad a_0=n(h-n+1),\quad u_0=a_0,\quad\kappa_0=h.
\]
Here $a-a_0=n(n-1)>0$. Thus the asymptotically sharp family in the previous
exterior cutoff was already reducible to an explicit middle state and a strictly
smaller endpoint E state. This does not lower that predecessor's bound on
\emph{all individual states}; it gives an existence-preserving way to remove
this branch from a search for genuinely E-only inputs.

The hard restriction is necessary. At $p=17$, the original E state
$(a,u,h,r,s,\kappa)=(5,5,5,1,1,6)$ has $\alpha=2,\beta=-2$, so its nonsingular
shape has radius two.

\section{The first possible nonsingular shape at radius eight}
\begin{theorem}[Exact next-profile reduction]\label{thm:eight}
At a hard prime, any nonsingular original E state with shape radius at most eight
must have
\[
 \boxed{(\alpha,\beta,h)=(8,-4,11),\qquad
 4r^2=11s^4-4s^2-3,\qquad p=44rs-11s^2+8.}               \tag{18}
\]
Here $r,s$ are odd and coprime. Necessary hard-class restrictions include
\[
 rs\equiv-1\pmod3,\qquad r^2\equiv s^2\equiv1\pmod5,
 \qquad rs\equiv1\pmod5.
\]
\end{theorem}
\begin{proof}
Only profiles containing $8$ or $-8$ are new beyond (16); their other coordinate
is in $\{-8,-4,4,8\}$. Any positive coordinate 4 is excluded by
Lemma~\ref{lem:char}. The products of absolute value 64 give $B=63$ or $-65$.
The first has only small prime factors. In the second, the only possible
nonresidue grades are 13 and 65, both $1\pmod4$, whereas $\alpha\equiv0$ and
$R\equiv3$ would require $hs^2\equiv3\pmod4$.

The pair $(-8,-4)$, in either order, has $B=31$ and forces $h=31,s$ odd.
Equation (17) gives $4r^2\equiv12\pmod{16}$, impossible.
For the remaining $B=-33$ pairs the only possible grades are 11 and 33.
The grade 33 is excluded modulo four. At grade 11, equation (17) and the prime
congruence modulo eight leave only $(8,-4)$; it makes $r,s$ odd and gives (18).
For completeness, the twelve new ordered pairs and every grade divisor
containing a prime greater than seven give exactly the following sixteen rows:
\[
\begin{array}{c|c|c}
(\alpha,\beta)&B&h\\ \hline
(-8,-4),(-4,-8)&31&31\\
(-8,4),(4,-8)&-33&11,33\\
(-8,8),(8,-8)&-65&13,65\\
(-4,8),(8,-4)&-33&11,33\\
(4,8),(8,4)&31&31.
\end{array}
\]
The omitted pairs $(-8,-8)$ and $(8,8)$ have $B=63$, supported only on
$3$ and $7$. Thus no admissible radius-eight grade row is hidden by the table.
This calculation can equivalently be exhausted over $r,s\pmod{24}$ using (17)
modulo 48, as in the certificate.

Modulo three, $3\mid s$ would force $3\mid r$, violating coprimality.
Thus $r,s$ are nonzero modulo three, and $p\equiv1$ gives $rs\equiv-1$.
Modulo five, (18) allows only $s^2=r^2=1$. Then $p\equiv2-rs\pmod5$;
the hard residue condition $p\equiv1,4\pmod5$ forces $rs\equiv1$.
Modulo seven, (18) gives
\[
 r^2\equiv s^4-s^2+1\pmod7.
\]
The square classes show that either $s\equiv0$, $r^2\equiv1$, $p\equiv1$,
or $s^2=r^2=1$; in the latter case the hard residue condition forces
$rs\equiv-1$ and $p\equiv2\pmod7$. Together with $p\equiv1\pmod5$ and
$p\equiv1\pmod{24}$ this sharpens the possible hard classes to
\[
 p\bmod840\in\{1,121\}.
\]
\end{proof}

\begin{corollary}[Radius eight already has an original middle coefficient]
Every hard-prime E state of shape radius at most eight forces an original M
state. In the singular case this is (15). In the sole possible nonsingular
case (18), the exact return is
\[
 \boxed{s_M=(p+3)/11,\quad
 (h_M,r_M,s_M,\lambda_M)=(1,3,(p+3)/11,1),}
\]
\[
 \boxed{a_M=3(p+3)/11,\quad u_M=9,\quad R_M=(p+36)/11,
 Q_M=11,\quad (x,y,z)=(a_M,ps_M,3p).}                    \tag{18a}
\]
It attains the preceding hard-prime middle localization equality
$11a_M=3(p+3)$.
\end{corollary}
\begin{proof}
The grade identity (6) gives $p\equiv\alpha=8\pmod{11}$, so $s_M$ is integral.
Then $p=11s_M-3=4\cdot3s_M-(3+s_M)$. A common factor of $3$ and $s_M$
would divide the hard prime $p$, impossible. The first-half inequalities follow
because $p\equiv1\pmod8$ gives $s_M\equiv4\pmod8$, hence $s_M>3$.
Substitution proves both the original M gate and its ordered
identity. The original E coordinates remain in the graph of this map; if they
are forgotten, its fibre is the set of positive coprime quartic pairs in (18)
with the fixed value of $p$. Theorem~\ref{thm:curve} bounds that fibre by three.
No integral-point classification is needed to prove this return.
\end{proof}

Equation (18) remains a specific integral-point problem for classifying the
exterior records, but no longer an obstruction to returning an original middle
state from any such record. The bounded scan in the certificate checks
odd $s\le200000$ exactly. No absence beyond that range, effective integral-point
cutoff, or global ES conclusion is inferred. The small point $r=s=1$ gives
$p=41$, outside the hard classes. The elementary Pell rewriting
\[
 (11s^2-2)^2-11(2r)^2=37
\]
also retains the extra condition that $(X+2)/11$ be a square; dropping that
condition would change the problem.

\section{A natural cubic involution fails the required fixed-prime return}
For $B\ne0$, (4) has the explicitly checked rational involution
\[
 \Psi(X,Y)=(B/X,BY/X^2),\qquad X\ne0.                    \tag{19}
\]
Substitution proves preservation of (4), and applying it twice is the identity.
For an original positive state it can have positive $X'$ only when $B>0$.
Integrality of $X'$ requires
\[
 v\mid B\iff v\mid4u\iff s^2\mid4r^2\iff s\mid2.
\]
If $s=2$, then $r$ is odd, $X'\equiv3\pmod4$, and $Y'=rX'$ is odd, so the
inverse $a'=Y'/2$ fails integrality. If $s=1$, the complete coordinate transport is
\[
 \begin{aligned}
 h'=v'&=4r^2+h-R-D, & r'&=r, & s'&=1,\\
 R'&=4r^2-D, & D'&=4r^2-R,\\
 a'&=rh', & u'&=h'r^2, & \kappa'&=D'-r,\\
 p'&=4rh'-R'.
 \end{aligned}                                          \tag{20}
\]
Whenever its signs are positive, it has $R',D'\equiv1\pmod4$ and
$p'\equiv3\pmod4$, not the original hard prime. More exactly,
\[
 \boxed{p'-p=(4r-1)(4r^2-R-D).}                          \tag{21}
\]
The second factor is $2\pmod4$, hence cannot vanish.
Thus (19) is an actual rational correspondence, but not a same-$p$ selector.
At $p=2129$, its positive composite-parameter image has
$p'=5951=11\cdot541$, $R'=D'=249$,
$h'=155,r'=10,s'=1$. This records what the map does rather than treating its
failure as evidence that the two geometries are unrelated.

The obvious reverse attempt from the always available residual-one source in
the class $p'\equiv3\pmod4$ also has an exact obstruction. Write
$p'=4hrs-1$, $R=1$, $D=4hr^2+1$, $v=hs^2$ with $\gcd(r,s)=1$.
For (19),
\[
 X'=4r^2/s^2+hs^2-4hr^2-2.
\]
An integer image still requires $s\mid2$. At $s=1$,
$X'=h-2-4(h-1)r^2<0$. At $s=2$,
$X'=4h-2-(4h-1)r^2\le-1$. For $s>2$ the image is not integral.
Thus the complete residual-one source does not supply a positive integral
inverse for the hard-prime problem. This does not prohibit other positive
$p'\equiv3\pmod4$ parameter records; the displayed image at $5951$ is one such
composite record. It identifies why the immediate reverse construction is
insufficient.

\section{Continuation status and executable boundary}
The legitimate new reductions are: the diagonal source has a same-prime middle
return with smaller cofactor; all hard-prime radius-seven E states are singular
and already have M hits at the same shell; the next radius has the unique
nonsingular quartic (18), whose grade congruence itself forces (18a). The marked curve includes the original fixed-$p$ line
(8), which its natural involution does not preserve.

None of this proves that an arbitrary prescribed prime supplies an original E
or M state. The universal Turn 7 existence milestone remains outstanding. The
finite profile classification is exhaustive because its domain consists of the
52 explicitly enumerated integer profiles; the bound on the quartic's search
prefix is a different, finite verification scope.

The supplied main verifier enumerates all original E square divisors at every
$p\equiv1\pmod8$ prime through its declared bound. The separate implementation
uses primitive $(h,r,s)$ enumeration and a different modular profile test. Both
keep the original integers and compare source fibres rather than comparing only
state totals. All check outcomes, failures and time limits are recorded by the
execution runner. No independent human mathematical review or Lean build is claimed.
